\section{Objetivos y justificación}
\subsection{Objetivo general}
Caracterizar el uso de herramientas digitales con fines académicos entre estudiantes universitarios de instituciones de educación superior de Bogotá, considerando los tipos de herramientas utilizadas, la frecuencia, las finalidades y las percepciones de utilidad y dificultad.

\subsection{Objetivos específicos}
\begin{multicols}{2}
    \begin{enumerate}
        \item Identificar las categorías de herramientas digitales que los estudiantes reportan utilizar para realizar actividades académicas.
        \item Describir la frecuencia de uso de las herramientas digitales y las actividades académicas para las que se emplean.
        \vspace{1cm}
        \item Describir la percepción de utilidad y las principales dificultades reportadas en el uso académico de estas herramientas.
        \item Explorar descriptivamente si las respuestas varían según características académicas seleccionadas, como nivel de formación, modalidad de estudio y semestre o periodo cursado, siempre que el tamaño y la composición de la muestra lo permitan.
    \end{enumerate}
\end{multicols}

\subsection{Justificación teórica}
El estudio permite aplicar conceptos metodológicos de investigación a un fenómeno educativo actual. La
UNESCO (2023) plantea que la presencia de tecnología no equivale automáticamente a una transformación
educativa; por eso es necesario analizar prácticas y condiciones de uso, no solo disponibilidad. La investigación
aportaría una descripción situada y podría servir como antecedente para preguntas posteriores más específicas.
\subsection{Justificación práctica}
Los resultados podrían orientar decisiones sobre inducción a plataformas institucionales, uso responsable
de recursos digitales, formación en búsqueda y evaluación de información, colaboración en línea y
acompañamiento a estudiantes que reporten dificultades. Las recomendaciones deberán derivarse de los datos
efectivamente obtenidos y no de expectativas previas del investigador.
\subsection{Justificación social y académica}
La investigación considera que el acceso, la disponibilidad de dispositivos, la conectividad y las
competencias digitales pueden ser desiguales. Aunque el DANE ofrece indicadores agregados de conectividad, el
estudio propuesto se concentra en prácticas universitarias que no quedan descritas por esas cifras generales. La
propuesta puede fomentar una discusión informada sobre inclusión y uso crítico de la tecnología.
\subsection{Viabilidad}
Una encuesta breve en línea o presencial puede aplicarse con recursos limitados, siempre que se obtenga
autorización para contactar participantes y se defina un procedimiento de muestreo transparente. Si no se logra
acceso a varias instituciones, el alcance deberá reducirse explícitamente a la institución o grupo accesible, sin
generalizar a toda Bogotá
